% 第 5 章：直流故障水平
\section{直流故障水平}\label{sec:sc-dc}

\subsection{电阻性刚性源模型}
\srcpath{dc\_bus\_fault\_level}（\srcpath{dc\_short\_circuit.hpp}）给出直流母线金属性故障的稳态
电流估计：把直流网络化简为其支路电阻图，电压参考（\texttt{DC\_V}）母线视作理想电压源。对
非源母线 $k$：
\begin{equation}
G_{NN}V_N=-G_{NS}V_S,\qquad
I_f=\frac{V_N[k]}{(G_{NN}^{-1})_{kk}+\max(0,R_f)},
\end{equation}
其中 $R_{kk}$ 为从母线 $k$ 到源的戴维南电阻（\fld{r\_thevenin\_pu}），$R_f$ 为故障路径电阻
（\fld{fault\_resistance\_pu}，默认 0 即金属性故障）。当分母不大于 $10^{-12}$、
$G_{NN}$ 不可逆或故障母线本身是 \texttt{DC\_V} 时返回未求解。标幺电流乘
$S_b/U_b$ 得 kA。依据为 \srcpath{src/short\_circuit/dc\_short\_circuit.cpp:dc\_bus\_fault\_level}。
这是\textbf{线路电阻限制}的准稳态故障水平，可用于直流断路器/电缆的第一层筛查；由于模型同时
省略可能抬高峰值的电容放电和可能压低稳态电流的换流器限流，它相对真实暂态\textbf{既不保证
上界也不保证下界}。

\subsection{直流断路器参与}
选项 \fld{DCFaultOptions} 控制断路器如何进入电导图：
\begin{itemize}
  \item \fld{consider\_dc\_breakers}：将 \texttt{DCCircuitBreaker} 纳入电导图；
  \item \fld{dc\_breakers\_control\_branches}：匹配的断路器控制对应 \texttt{DCBranch}——断开阻断支路、
        闭合串入接触电阻；
  \item \fld{add\_unassigned\_closed\_breaker\_edges}：未指派的闭合断路器作独立导电边；
  \item \fld{min\_closed\_breaker\_resistance\_pu}：正电阻闭合边的数值下限；严格零电阻边作为
        理想导体先由并查集收缩，不以小电阻替代拓扑等势关系。
\end{itemize}
逐断路器故障责任记于 \fld{DCBreakerDutyResult}：\fld{i\_duty\_ka}、\fld{i\_breaking\_ka}、
\fld{breaking\_rating\_ok}、\fld{controls\_branch} 与 \fld{r\_pu}。实现先以补偿定理恢复故障后
母线电压，再按每条实际电阻边 $I_{ij}=(V_i-V_j)/R_{ij}$ 计算责任；串联 DCCB 电流相同，并联
路径按电阻分流，开路支线为 0。端点相同且无法唯一归属到并行支路的 DCCB 显式拒绝，不猜测。
结果 \fld{DCFaultResult} 另记
\fld{v\_prefault\_pu}、边计数（\fld{dc\_branch\_edges\_used}/\fld{dccb\_edges\_used}）与
\fld{dccb\_open\_count}/\fld{dccb\_blocked\_branch\_count}。

\subsection{边界}
本估计\textbf{不}建模换流器限流、电容放电暂态、线路电感、源控制与直流保护跳闸。
真实峰值放电电流可能更高（电容），而稳态受控电流可能更低（换流器限流），故不能视为真实
故障波形的包络。需要暂态峰值、$di/dt$ 或
受控稳态时应改用暂态仿真（\srcpath{docs/theory/transient\_runtime.md}）。
